In a Boolean ring, $x=x^2$. More generally, $x^n=x$ with $n>1$ gives $x=x^{n-2}x^2$. In both cases $(x)=(x^2)$, so \exref{2.27} gives absolute flatness.

If $A$ is absolutely flat, write $x=ax^2$ for each $x\in A$. This identity descends to every quotient, which is therefore absolutely flat by \exref{2.27}.

Suppose $A$ is local and absolutely flat. If $x$ is a nonunit, write $x=ax^2$. Then $e=ax$ is idempotent and lies in the maximal ideal, so $1-e$ is a unit. Since $e(1-e)=0$, we get $e=0$ and $x=xe=0$. Thus every nonzero element is a unit, and $A$ is a field.

In any absolutely flat ring, $x=ax^2$ gives $x(1-ax)=0$. If $x$ is a nonunit, then $1-ax\neq0$, so $x$ is a zero-divisor.
